\documentclass[a4paper,14pt]{article}
\usepackage[spanish]{babel}
\usepackage{pgfornament}
\usepackage{tikz}
\usetikzlibrary{patterns,patterns.meta}
\usepackage[T1]{fontenc}
\usepackage[sfdefault]{ClearSans}
%210 x 297 
%52.5x74
\newcommand{\recibon}{00001}
\newcommand{\persona}{Ferreyra Gustavo David}
\newcommand{\pesos}{Mil quinientos setenta y siente}
\newcommand{\pesosn}{1577}
\newcommand {\concepto}{REBOQUE FINO}

\begin{document}

\begin{tikzpicture}[overlay,shift={(-5,-2.5)}]
\draw [dashed](0,0) rectangle (21,7.4);
\draw  (3,6) node [right]{\Large N$^\circ$: \recibon};
\draw (20,6) node [left]{\Large San Rafael Mendoza, \today  };
\draw (3,5) node [right] {\Large Recibí de \persona.};
\draw (3,4) node [right] {\Large La cantidad de pesos: \pesos.-};
\draw (3,3) node [right] {\Large En concepto de: \concepto.-};
\draw (3,2) node [right] {\Large Son: \$\pesosn.-};
\draw [rounded corners=0.5cm](11,1) rectangle+(9,1.5);
\draw [rounded corners,pattern={Dots[angle=45,distance={3pt}]}](4.2,1.5) rectangle (8,2.5);
\draw [rounded corners,pattern={Dots[angle=45,distance={3pt}]}](8,3.5) rectangle +(12,1);
\draw [line width=1mm,rounded corners=1cm](.5,.5)rectangle(20.5,7);
\draw (15.5,1.2) node {Firma};
\draw (2,5.5) node {\pgfornament[color=gray,scale=0.5,symmetry=v]{26}};
\draw (2,4) node {\pgfornament[color=gray,scale=0.5]{25}};
\draw (2,2.3) node {\pgfornament[color=gray,scale=0.5,symmetry=h]{26}};

\end{tikzpicture}
\end{document}